\documentclass[a4paper]{article}
\usepackage{amsfonts}
\usepackage{amsmath}
\usepackage{amssymb}
\usepackage{graphicx}     

\begin{document}
  
\noindent \large Auteur: Anna Voronovska \\
\noindent \large Studierichting: Informatica\\
\noindent \large Oefeningenreeks 7E nr. 1.6\\

\medskip

\normalsize

$f(x, y)= x^3+y^3-6xy$\\

$\dfrac{\partial f}{\partial x} = 3x^2 -6y$ (*)\\

$\dfrac{\partial f}{\partial y} = 3y^2 -6x$ (**)\\

Kritieke punten:\\

(*) $3x^2-6y=0$\\

    $x^2=2y$\\

    $y=\dfrac{x^2}{2}$\\
    
(**) $3y^2 -6x=0$\\

	 $y^2=2x$\\
	 
Wij substitueren $y=\dfrac{x^2}{2}$ in $y^2=2x$:\\

$\left(\dfrac{x^2}{2}\right)^2 =2x$\\

$\dfrac{x^4}{4}=2x$\\

$x^4-8x=0$\\

$x(x^3-8)=0$\\

$x_1=0$ \\ 

$x^3-8=0$\\

$x^3=8$\\

$x_2=2$\\

Voor $x=0$: $y = \dfrac{0^2}{2} = 0$\\

Voor $x=2$: $y = \dfrac{2^2}{2} = 2$\\

Kritieke punten: $(0,0)$ en $(2,2)$.\\

$H = 
\left|
\begin{array}{ccc}
\dfrac{\partial^2 f}{\partial x^2} & \dfrac{\partial^2 f}{\partial x \partial y}\\
\dfrac{\partial^2 f}{\partial x \partial y} & \dfrac{\partial^2 f}{\partial y^2}
\end{array}
\right| = 
\left|
\begin{array}{ccc}
6x & -6\\
-6 & 6y
\end{array}
\right| = (6x)(6y)-(-6)^2=36xy-36
$\\




Als $(0,0)$: $H = 36\cdot 0 \cdot 0 -36 = -36$\\

Dus $(0, 0)$ is een zadelpunt.\\

Als $(2, 2)$: $H = 36 \cdot 2 \cdot 2 - 36 = 144 - 36 = 108$\\

Dus $(2, 2)$ is een lokaal minimum in:\\

$f(2, 2) = 2^3 + 2^3 -6 \cdot 2 \cdot 2 = 8+8-24=-8$\\
\begin{thebibliography}{999}
%\bibitem{a} Referentie
\end{thebibliography}
\end{document} 
